\documentclass[11pt]{article}

\usepackage[margin=1in]{geometry}
\usepackage{amsmath,amssymb,amsthm,mathtools}
\usepackage{hyperref}
\usepackage{enumitem}

\title{Aspectual Counting Framework (ACF):\\
Aspects, Contexts, Bridges, and Higher-Order Conflict Geometry}
\author{Ryan O. Van Gelder}
\date{December 29, 2025}

% ------------------ Theorem environments ------------------
\theoremstyle{definition}
\newtheorem{definition}{Definition}[section]
\newtheorem{axiom}{Axiom}[section]
\newtheorem{assumption}{Assumption}[section]

\theoremstyle{plain}
\newtheorem{theorem}{Theorem}[section]
\newtheorem{proposition}{Proposition}[section]
\newtheorem{lemma}{Lemma}[section]
\newtheorem{corollary}{Corollary}[section]

\theoremstyle{remark}
\newtheorem{remark}{Remark}[section]

% ------------------ Notation helpers ------------------
\newcommand{\Pow}{\mathcal{P}}
\newcommand{\Obj}{\mathrm{Obj}}
\newcommand{\FinSet}{\mathbf{FinSet}}
\newcommand{\bbN}{\mathbb{N}}
\newcommand{\bbR}{\mathbb{R}}
\newcommand{\EqCl}{\mathrm{EqCl}}
\newcommand{\argmax}{\operatorname*{arg\,max}}
\newcommand{\OPT}{\mathrm{OPT}}
\newcommand{\Sat}{\mathrm{Sat}}
\newcommand{\Unsat}{\mathrm{Unsat}}
\newcommand{\MUS}{\mathrm{MUS}}
\newcommand{\load}{\mathrm{load}}
\newcommand{\Feas}{\mathrm{Feas}}
\newcommand{\nuFrac}{\nu_f}

\begin{document}
\maketitle

\begin{abstract}
The Aspectual Counting Framework (ACF) formalizes the idea that numerical answers depend on the \emph{aspect} under which we count. An aspect fixes (i) admissible units, (ii) invariances, and (iii) permissible inferences within and across aspects via explicit \emph{bridges}. This article develops ACF from a set-based quotient model through a functorial formulation, introduces contexts as a semilattice of practices, defines a bridge-selection principle as a prioritized satisfiable-subset problem, and analyzes the resulting conflict geometry via graphs, hypergraphs, Helly-type conditions, and a $k$-shadow hierarchy. We provide a concrete inconsistency theorem for Ship-of-Theseus-style paradoxes (as a non-existence of a global identity equivalence satisfying incompatible bridges), show how selection reduces to maximum-weight independent set (graph case) or its hypergraph analogue (higher-order case), define a practical Helly order induced by weights, give local dominance and antidominance criteria, and present a sound hypergraph-native certification procedure with a provable runtime-bounded local candidate generator. The only heuristic component is which candidate sets are tested; every certified constraint is sound.
\end{abstract}

\tableofcontents

\section{Motivation: aspectual answers and Theseus}
In counting disputes (Ship of Theseus, restored artifacts, legal identity, software versions), parties often share all base-level facts yet disagree numerically. ACF posits: such disagreements frequently arise from \emph{unlicensed cross-aspect inference}. One can count planks, ship-continuants, replacement events, legal registrations, functional roles, etc.; these are different invariants of the same underlying situation.

\section{ACF-1: units as quotients and invariance}

\subsection{Situations and aspects (quotient form)}
Let $E$ be a domain of items (tokens, time-slices, events, etc.). Let $\Obj$ be a class of situations $X$ with carriers $|X|\subseteq E$.

\begin{definition}[Aspect as unitization]\label{def:unitization}
An \emph{aspect} $\alpha$ assigns to each $X\in\Obj$:
\begin{enumerate}[label=(\roman*)]
\item a selection $S_\alpha(X)\subseteq |X|$ (eligible candidates),
\item an equivalence relation $\sim_{\alpha,X}$ on $S_\alpha(X)$ (when candidates count as the same unit).
\end{enumerate}
Define the set of $\alpha$-units $U_\alpha(X):=S_\alpha(X)/\sim_{\alpha,X}$.
\end{definition}

\begin{definition}[Aspectual count]
If $U_\alpha(X)$ is finite, define
\[
a(X,\alpha):=|U_\alpha(X)|\in\bbN.
\]
\end{definition}

\subsection{Invariance by construction}
\begin{definition}[$\alpha$-preserving bijection]
A bijection $f:|X|\to |Y|$ is \emph{$\alpha$-preserving} if
\begin{enumerate}[label=(\roman*)]
\item $x\in S_\alpha(X)\iff f(x)\in S_\alpha(Y)$,
\item $x\sim_{\alpha,X}x'\iff f(x)\sim_{\alpha,Y}f(x')$.
\end{enumerate}
\end{definition}

\begin{theorem}[Invariance]
If $f$ is $\alpha$-preserving, then $a(X,\alpha)=a(Y,\alpha)$.
\end{theorem}

\subsection{Locality and additivity}
\begin{axiom}[Restriction coherence]\label{ax:restriction}
If $Y\subseteq |X|$ and $X\!\upharpoonright\!Y$ is defined, then
$S_\alpha(X\!\upharpoonright\!Y)=S_\alpha(X)\cap Y$ and $\sim_{\alpha,X\upharpoonright Y}$ is the restriction of $\sim_{\alpha,X}$.
\end{axiom}

\begin{definition}[$\alpha$-separation]
If a disjoint sum $X\sqcup Y$ is defined, say $X\perp_\alpha Y$ if
$U_\alpha(X\sqcup Y)\cong U_\alpha(X)\sqcup U_\alpha(Y)$.
\end{definition}

\begin{axiom}[Conditional additivity]\label{ax:add}
If $X\perp_\alpha Y$ then $a(X\sqcup Y,\alpha)=a(X,\alpha)+a(Y,\alpha)$.
\end{axiom}

\begin{remark}
Continuant aspects often violate additivity over temporal concatenation: two history segments each contain one continuant, yet their concatenation also contains one.
\end{remark}

\section{ACF-2: functorial aspects, contexts, and bridges}

\subsection{Aspects as functors}
Let $\mathcal O$ be a category of situations (objects are situations, morphisms are structure-preserving maps; isomorphisms are renamings).

\begin{definition}[Aspect functor]\label{def:functor}
An aspect is a functor $U_\alpha:\mathcal O\to\FinSet$.
Define $a(X,\alpha):=|U_\alpha(X)|$ when finite.
\end{definition}

\begin{theorem}[Functorial invariance]
If $g:X\cong Y$ in $\mathcal O$, then $U_\alpha(g):U_\alpha(X)\cong U_\alpha(Y)$ and $a(X,\alpha)=a(Y,\alpha)$.
\end{theorem}

\begin{remark}
The quotient model of Definition~\ref{def:unitization} is recovered by taking $U_\alpha(X)=S_\alpha(X)/\sim_{\alpha,X}$.
\end{remark}

\subsection{Contexts as admissibility constraints}
A key philosophical point is that admissibility depends on what a practice can \emph{talk about}.

\begin{definition}[Context-indexed situation categories]
A context $c$ supplies a category $\mathcal O_c$ of situations described in a language/signature $L_c$, and a forgetful/translation functor $F_c:\mathcal O\to\mathcal O_c$.
\end{definition}

\begin{definition}[Context-admissible aspect]
An aspect is \emph{admissible in context $c$} if there exists a functor
$\overline U_{\alpha,c}:\mathcal O_c\to\FinSet$ such that
$U_\alpha\cong \overline U_{\alpha,c}\circ F_c$.
\end{definition}

\subsection{Refinement as natural surjection}
\begin{definition}[Refinement]
An aspect $\beta$ \emph{refines} $\alpha$ (write $\beta\succeq\alpha$) if there exists a natural transformation
$\rho:U_\beta\Rightarrow U_\alpha$ such that each component $\rho_X:U_\beta(X)\twoheadrightarrow U_\alpha(X)$ is surjective.
\end{definition}

\begin{proposition}[Monotonicity under refinement]
If $\beta\succeq\alpha$ and both counts are finite, then $a(X,\beta)\ge a(X,\alpha)$.
\end{proposition}

\subsection{Bridges as natural relations}
\begin{definition}[Bridge as natural relation]\label{def:bridge}
Fix a context $c$. A bridge $\tau:U_\alpha \rightsquigarrow U_\beta$ is a natural family of relations
$\tau_X\subseteq U_\alpha(X)\times U_\beta(X)$.
\end{definition}

\begin{remark}
Bridges are often partial, context-dependent, and relational rather than functional. ACF enforces: \emph{no bridge, no transfer}.
\end{remark}

\section{Ship of Theseus as non-existence of a global identity}
\subsection{A minimal inconsistency model}
Let $S=\{s_0,s_1,s_2\}$ be three ship-slices: original at $t_0$, maintained at $t_1$, reassembled at $t_1$.
Let $\#\subseteq S^2$ encode co-present distinctness: $s_1\# s_2$.
Let $C$ (continuity) relate $s_0$ to $s_1$, and $M$ (material provenance) relate $s_0$ to $s_2$.

\begin{theorem}[Theseus inconsistency theorem]\label{thm:theseus}
There is no equivalence relation $\sim$ on $S$ such that:
(i) $C(x,y)\Rightarrow x\sim y$,
(ii) $M(x,y)\Rightarrow x\sim y$,
(iii) $x\# y\Rightarrow \neg(x\sim y)$.
\end{theorem}

\begin{proof}
From (i) and $C(s_0,s_1)$, $s_0\sim s_1$.
From (ii) and $M(s_0,s_2)$, $s_0\sim s_2$.
Transitivity yields $s_1\sim s_2$, contradicting (iii) since $s_1\# s_2$.
\end{proof}

\begin{remark}
The ``paradox'' is the implicit demand that one identity relation satisfy incompatible bridge-induced constraints.
\end{remark}

\section{ACF-3: contexts as a semilattice and bridge selection}

\subsection{Contexts as a join-semilattice}
\begin{definition}[Context semilattice]
Let $(\mathcal C,\le)$ be a poset of contexts. Assume a join operation $c\vee d$ (least context extending both). Intuitively: $c\vee d$ combines practices, constraints, and bridge libraries.
\end{definition}

\subsection{$\#$-strict contexts and satisfiability}
Fix $c\in\mathcal C$.
Let $B_c$ be available bridge schemata. Each schema $b\in B_c$ induces for each situation $X$ a required-identification relation $R_{b,c}(X)$ on candidate slices $S_c(X)$.

Define forced identity for $\Gamma\subseteq B_c$ on $X$ by
\[
\equiv_{\Gamma,c}(X):=\EqCl\Bigl(\bigcup_{b\in\Gamma} R_{b,c}(X)\Bigr).
\]
Let $\#_c(X)\subseteq S_c(X)^2$ be co-presence distinctness.

\begin{definition}[$\#$-strictness and satisfiability]
Context $c$ is \emph{$\#$-strict} if co-presence distinctness is non-negotiable.
A set $\Gamma\subseteq B_c$ is \emph{satisfiable} in $c$ if for all $X$,
\[
\equiv_{\Gamma,c}(X)\cap \#_c(X)=\varnothing.
\]
\end{definition}

\subsection{Prioritized satisfiable-subset selection}
Let $\preceq_c$ be a priority preorder on $B_c$, inducing levels $L_1,\dots,L_K$ (highest to lowest).
Encode lexicographic priorities by weights $w(b)=M^{K-i}$ for $b\in L_i$, $M>|B_c|$.

\begin{definition}[Selection problem]
Choose $\Gamma^\*\subseteq B_c$ maximizing $\sum_{b\in\Gamma} w(b)$ subject to $\Gamma$ satisfiable.
\end{definition}

\section{Conflict geometry: from graphs to hypergraphs}

\subsection{Monotonicity and minimal unsatisfiable sets}
\begin{lemma}[Monotonicity]
If $\Gamma\subseteq\Delta$ and $\Unsat(\Gamma)$ then $\Unsat(\Delta)$.
\end{lemma}

\begin{definition}[Minimal unsatisfiable sets (MUS)]
A set $M\subseteq B_c$ is a MUS if $\Unsat(M)$ and every proper subset is satisfiable.
Let $\mathcal M_c$ denote the family of all MUS.
\end{definition}

\begin{lemma}[MUS basis]
$\Unsat(\Gamma)$ iff $\exists M\in\mathcal M_c$ with $M\subseteq\Gamma$.
\end{lemma}

\subsection{Graph reduction axiom (BW2 / MUS2)}
\begin{axiom}[MUS2]\label{ax:mus2}
All minimal unsatisfiable sets have size $2$: $\forall M\in\mathcal M_c,\ |M|=2$.
\end{axiom}

Under MUS2, satisfiable sets are exactly independent sets in a graph.

\subsection{Conflict hypergraph (general case)}
\begin{definition}[Conflict hypergraph]
Define $H_c=(V_c,\mathcal E_c)$ with $V_c=B_c$ and $\mathcal E_c=\mathcal M_c$.
A bridge set $\Gamma$ is satisfiable iff it contains no hyperedge:
$\forall e\in\mathcal E_c,\ e\not\subseteq\Gamma$.
\end{definition}

\subsection{Reduction to (hyper)graph independent set}
\begin{proposition}[Optimization form]
Bridge selection is equivalent to:
\[
\max\Bigl\{\sum_{b\in\Gamma} w(b): \Gamma\subseteq B_c,\ \forall M\in\mathcal M_c,\ \Gamma\not\supseteq M\Bigr\}.
\]
As an ILP:
\[
\max \sum_{b} w(b)x_b \quad \text{s.t.}\quad \sum_{b\in M} x_b \le |M|-1\ \ (\forall M\in\mathcal M_c),\ \ x_b\in\{0,1\}.
\]
\end{proposition}

\section{Helly-type characterization and the $k$-shadow hierarchy}

\subsection{Helly number of satisfiability}
\begin{definition}[Helly number]
Let $h(c):=\max\{|M|: M\in\mathcal M_c\}\in\bbN\cup\{\infty\}$.
\end{definition}

\begin{theorem}[Helly/MUS equivalence]\label{thm:helly}
For $k\ge 2$, the following are equivalent:
(i) every MUS has size $\le k$;
(ii) satisfiability is $k$-Helly (if all subsets of size $\le k$ are satisfiable then the whole set is satisfiable);
(iii) $\Unsat(\Gamma)\Rightarrow \exists \Delta\subseteq\Gamma$ with $|\Delta|\le k$ and $\Unsat(\Delta)$.
\end{theorem}

\begin{corollary}[Graph reduction]
The conflict hypergraph reduces to a graph iff $h(c)=2$ (equivalently MUS2).
\end{corollary}

\subsection{$k$-shadow approximation}
\begin{definition}[$k$-shadow of MUS family]
For a MUS $M$ and $k\ge 2$, define
\[
\partial_k(M):=
\begin{cases}
\{S\subseteq M:|S|=k\}, & |M|>k,\\
\{M\}, & |M|\le k.
\end{cases}
\qquad
\partial_k\mathcal M_c:=\bigcup_{M\in\mathcal M_c}\partial_k(M).
\]
Define the $k$-shadow approximation hypergraph $H^{(k)}=(B_c,\partial_k\mathcal M_c)$.
\end{definition}

\begin{proposition}[Equivalent per-MUS cap]
$\Gamma$ is feasible in $H^{(k)}$ iff for all $M\in\mathcal M_c$,
\[
|\Gamma\cap M|\le \min(k-1,|M|-1).
\]
\end{proposition}

\begin{theorem}[Monotone interpolation]
Let $OPT_k$ be the maximum weight over $\Feas_k$ (the feasible sets of $H^{(k)}$), and $OPT_\star$ the true optimum.
Then
\[
OPT_2 \le OPT_3 \le \cdots \le OPT_{h(c)} = OPT_\star.
\]
\end{theorem}

\section{Practical Helly order and early exactness under weights}

\subsection{Practical Helly order}
\begin{definition}[Practical Helly order]
Define
\[
h_w(c):=\min\{k\ge 2:\ OPT_k=OPT_\star\}.
\]
\end{definition}

\begin{definition}[MUS-load]
For $\Gamma\subseteq B_c$, define $\load(\Gamma):=\max_{M\in\mathcal M_c}|\Gamma\cap M|$.
Let
\[
L^\* := \min\{\load(\Gamma): \Gamma\in\argmax_{\Delta\in\Feas_\star} w(\Delta)\}.
\]
\end{definition}

\begin{theorem}[Exactness criterion via optimal-load]\label{thm:load}
For $k\ge 2$,
\[
OPT_k = OPT_\star \quad\Longleftrightarrow\quad k\ge L^\*+1.
\]
Consequently $h_w(c)=L^\*+1$ and $2\le h_w(c)\le h(c)$.
\end{theorem}

\section{Local dominance: sufficient conditions for early exactness}

\subsection{Hypergraph-native dominance (MUS-incidence)}
Let $\mathcal M(b):=\{M\in\mathcal M_c: b\in M\}$.

\begin{definition}[Incidence dominance]
Define $d\succeq b$ if $w(d)\ge w(b)$ and $\mathcal M(d)\subseteq \mathcal M(b)$.
\end{definition}

\begin{lemma}[Safe swap]
If $d\succeq b$ and $\Gamma$ is satisfiable with $b\in\Gamma$, then $(\Gamma\setminus\{b\})\cup\{d\}$ is satisfiable and has weight no smaller.
\end{lemma}

\begin{definition}[$k$-tail]
For each MUS $M$, let $Top_{k-1}(M)$ be the $\min(k-1,|M|)$ highest-weight elements of $M$.
Define $Tail_k(M):=M\setminus Top_{k-1}(M)$ and $Tail_k:=\bigcup_M Tail_k(M)$.
\end{definition}

\begin{assumption}[Local dominance condition $DC(k)$]
For every $b\in Tail_k$, there exists $d\neq b$ with $d\succeq b$.
\end{assumption}

\begin{theorem}[Local dominance implies $h_w(c)\le k]
If $DC(k)$ holds in a $\#$-strict context, then $OPT_k=OPT_\star$, hence $h_w(c)\le k$.
\end{theorem}

\subsection{Graph-approx dominance (2-section neighborhood)}
Define the 2-section graph $G^{2sec}$ on $B_c$ where $\{u,v\}$ is an edge iff they co-occur in some MUS.
Let $N[b]$ denote closed neighborhood in $G^{2sec}$.

\begin{definition}[Graph dominance]
$d\succeq_G b$ if $w(d)\ge w(b)$ and $N[d]\subseteq N[b]$.
\end{definition}

\begin{assumption}[Graph dominance condition $DC_G(k)$]
For every $b\in Tail_k$, there exists $d\neq b$ with $d\succeq_G b$.
\end{assumption}

\begin{theorem}[Tail elimination for MWIS on $G^{2sec}$]
If $DC_G(k)$ holds, then $G^{2sec}$ has a maximum-weight independent set that avoids $Tail_k$.
\end{theorem}

\section{Dual local lower bounds: antidomination structures}

\subsection{Local forcing of vertices (graph setting)}
Let $G$ be a weighted graph. For $b\in V(G)$, let $UB(b)$ be any upper bound on the maximum weight of an independent set in $G[N(b)\setminus\{b\}]$.

\begin{lemma}[Local forcing]
If $w(b)>UB(b)$, then every MWIS contains $b$.
\end{lemma}

\subsection{Weakened packing via an auxiliary graph and matching}
Fix $k$ and tail set $T=Tail_k$ in $G^{2sec}$. For $u,v\in T$ nonadjacent, let $Z(u,v)=N[u]\cup N[v]$.
If $w(u)+w(v)$ exceeds an upper bound on the best independent-set weight within $Z(u,v)\setminus\{u,v\}$, then every MWIS must include at least one of $\{u,v\}$.

Define an auxiliary graph $A_k$ on vertex set $T$, with edge $\{u,v\}$ if the above certificate holds.
Then any MWIS must hit each edge of a matching; hence matching size lower-bounds unavoidable tails.

\begin{theorem}[Matching lower bound]
Let $\mu_k$ be the maximum matching size in $A_k$. Then every MWIS $I^\*$ satisfies $|I^\*\cap Tail_k|\ge \mu_k$.
\end{theorem}

\section{Hypergraph-native certified unavoidable constraints and $r$-uniform packing}

\subsection{Certified unavoidable tail-hyperedges}
Fix tail level $k$ and choose $Q\subseteq Tail_k$ with $|Q|=r$.

Define
\[
\mathcal M(Q):=\{M\in\mathcal M_c: M\cap Q\neq\varnothing\},\qquad
Z(Q):=\bigcup_{M\in\mathcal M(Q)} (M\setminus Q),\qquad
\mathcal M_Z(Q):=\{M\in\mathcal M_c: M\subseteq Z(Q)\}.
\]

\subsection{Standard local bound: LP relaxation}
\begin{definition}[Local LP upper bound]
Define $UB_{\mathrm{LP}}(Q)$ as the optimum of
\[
\begin{aligned}
\max \ & \sum_{b\in Z(Q)} w(b)x_b\\
\text{s.t. } & \sum_{b\in M} x_b \le |M|-1 \quad(\forall M\in\mathcal M_Z(Q)),\\
& 0\le x_b \le 1 \quad(\forall b\in Z(Q)).
\end{aligned}
\]
\end{definition}

\begin{definition}[Certified unavoidable set]
Assuming $Q$ itself contains no MUS (i.e.\ $\forall M\in\mathcal M_c,\ M\nsubseteq Q$), call $Q$ \emph{certified unavoidable} if
\[
w(Q) > UB_{\mathrm{LP}}(Q).
\]
\end{definition}

\begin{theorem}[Unavoidability soundness]
If $Q$ is certified unavoidable, then every maximum-weight feasible selection $\Gamma^\*$ intersects $Q$.
\end{theorem}

\subsection{Auxiliary hypergraph and fractional matching lower bound}
Let $\mathcal E_{k,r}$ be the family of certified unavoidable $r$-sets.
Define $A_{k,r}=(Tail_k,\mathcal E_{k,r})$.

\begin{theorem}[Fractional matching tail lower bound]
Let $\nuFrac(A_{k,r})$ be the fractional matching number of $A_{k,r}$.
Then for every optimal selection $\Gamma^\*$,
\[
|\Gamma^\*\cap Tail_k| \ge \tau(A_{k,r}) \ge \nuFrac(A_{k,r}),
\]
where $\tau$ is the transversal number.
\end{theorem}

\begin{remark}
For $r=2$, this recovers the graph/matching bound in polynomial time. For $r\ge 3$, integral hypergraph matching is hard in general, but $\nuFrac$ remains polynomial via LP.
\end{remark}

\section{Algorithm specification and soundness guarantees}

\subsection{ACF certification procedure (compact)}
\paragraph{Inputs.}
$B$, weights $w$, MUS family $\mathcal M_c$, tail rule $Tail_k$, target $(k,r)$ values, and a candidate generator producing finite families $\mathcal Q_{k,r}\subseteq \binom{Tail_k}{r}$.

\paragraph{Outputs.}
Certified families $\mathcal E_{k,r}\subseteq \mathcal Q_{k,r}$, auxiliary hypergraphs $A_{k,r}$, and lower bounds $L_{k,r}:=\nuFrac(A_{k,r})$.

\paragraph{Certification.}
For each $Q\in\mathcal Q_{k,r}$:
(i) discard if it contains a MUS;
(ii) compute $Z(Q)$ and $\mathcal M_Z(Q)$;
(iii) solve the local LP to get $UB_{\mathrm{LP}}(Q)$;
(iv) if $w(Q)>UB_{\mathrm{LP}}(Q)$, add $Q$ to $\mathcal E_{k,r}$.

\paragraph{Bounding.}
Compute $\nuFrac(A_{k,r})$ via the fractional matching LP.

\subsection{Soundness theorem (certified vs heuristic)}
\begin{theorem}[Procedure soundness]\label{thm:soundness}
Fix a $\#$-strict context with MUS family $\mathcal M_c$ and weights $w$.
Run the certification procedure with any candidate sets $\mathcal Q_{k,r}$ (possibly generated heuristically).
Then:
\begin{enumerate}[label=(\roman*)]
\item Every certified $Q\in\mathcal E_{k,r}$ is a valid unavoidable constraint:
$\forall \Gamma^\*\in\OPT,\ \Gamma^\*\cap Q\neq\varnothing$.
\item Consequently every optimum hits every hyperedge in $A_{k,r}$.
\item The computed bound $L_{k,r}=\nuFrac(A_{k,r})$ satisfies
$\forall \Gamma^\*\in\OPT,\ |\Gamma^\*\cap Tail_k|\ge L_{k,r}$.
\item The only heuristic element is coverage: failing to test some $Q$ can miss constraints (false negatives), but the procedure cannot certify a false unavoidable constraint (no false positives), assuming LPs are solved correctly.
\end{enumerate}
\end{theorem}

\section{A minimal local candidate generator with runtime bound}

\subsection{MUS co-incidence neighborhood generator (MCN)}
Fix $k$ and $T=Tail_k$.

\paragraph{Co-incidence graph.}
Build $C_k=(T,E_k)$ with edge $\{u,v\}$ iff $\exists M\in\mathcal M_c$ such that $\{u,v\}\subseteq M$.
This can be built in time $O(\sum_{M\in\mathcal M_c} |M\cap T|^2)$.

\paragraph{Caps.}
Choose a neighborhood cap $s\in\bbN$ and radius $\ell\in\{1,2\}$.
For each $t\in T$, let $N^{(\ell)}(t)$ be its $\ell$-hop neighborhood in $C_k$.
Define $N_s(t)$ as the top $\min(s,|N^{(\ell)}(t)|)$ vertices in $N^{(\ell)}(t)$ by weight.

\paragraph{Star enumeration.}
For each target size $r$ and each $t\in T$, generate candidates
\[
\mathcal Q_{k,r}(t):=\{\{t\}\cup S: S\subseteq N_s(t),\ |S|=r-1\}.
\]
Let $\mathcal Q_{k,r}=\bigcup_{t\in T}\mathcal Q_{k,r}(t)$ (deduplicated).

\begin{proposition}[Generator runtime bound]
For fixed $k$, the MCN generator runs in
\[
O\!\left(\sum_{M\in\mathcal M_c} |M\cap T|^2\right)
\;+\;
O\!\left(|T|\sum_{r\in R}\binom{s}{r-1}\right)
\]
time (plus linear hashing for deduplication).
\end{proposition}

\begin{remark}
Soundness remains unconditional (Theorem~\ref{thm:soundness}); the generator only affects completeness.
\end{remark}

\section{Empirical and computational predictions}
\subsection{Graph vs hypergraph as a testable hypothesis}
The MUS size profile $p_k=\frac{|\{M\in\mathcal M_c:|M|=k\}|}{|\mathcal M_c|}$ and Helly number $h(c)$ quantify higher-order conflict.
Domains with $h(c)=2$ are effectively pairwise; domains with frequent $|M|\ge 3$ support ``all pairs ok, some triples forbidden'' phenomena.

\subsection{$k$-shadow saturation and order of interaction}
Compute $OPT_k$ (or bounds thereof) across $k$ to estimate $h_w(c)$.
Compute certified lower bounds on unavoidable tails via $L_{k,r}$; growth of $L_{k,r}$ with $r$ indicates higher-order interaction beyond pairs.

\subsection{Default parameter heuristics}
Typical regimes suggest starting with $r\in\{2,3\}$ and modest neighborhood caps (e.g.\ $s\approx 15$) with $\ell=1$, escalating to $\ell=2$ when MUS are larger or co-incidence neighborhoods are sparse.

\section{Conclusion and open directions}
ACF transforms counting/identity disputes into explicit choices of (i) aspects, (ii) contexts, and (iii) bridges. The resulting selection problem has a rich geometry: graph-like under MUS2, hypergraph-like otherwise, with Helly numbers measuring interaction order. The $k$-shadow hierarchy interpolates between pairwise approximations and true constraints, while $h_w(c)$ captures the effective order needed under a weight regime. Local dominance and certified unavoidability procedures yield fast, interpretable sufficient conditions and lower bounds. Future work includes: learning $\mathcal M_c$ from data, context morphisms and joins as empirical shifts, and characterizing when higher-order structure is necessary for stable disagreement.

\nocite{*}
\bibliographystyle{unsrt}
\bibliography{references}
\clearpage
\end{document}
